\documentclass{ximera}[font=12pt]

%Next 2 lines create spaces between items in itemize and enumerate
\setitemize{itemsep=2ex}
\setenumerate{itemsep=1.5ex}

%Creates space between paragraphs - Maybe???
\setlength{\parskip}{\medskipamount}

\title{Adding Integers: Financial Reasoning}   % Use xmlatex all -s to compile when SERVE refuses to work.
\author{Kelly Stady}
\license{CC: 0}         % replace with an appropriate license, or set it in xmPreamble

\begin{document}

\begin{abstract}

\end{abstract}

\maketitle
\label{xim:AddIntMoneyActivity}


\section*{Adding Integers Using Financial Reasoning}

Number lines and the Orange-Blue Rule are great for visualizing addition and finding answers to simple problems, but they aren't 
always practical.  For example, we would need a very large number line to find the value of $-185 + 99.$  Instead, we can add 
integers by thinking in terms of money.


\begin{itemize}
\item \textbf{Assets (+)} include money that you \textbf{have}, a gift you receive or a deposit.

\item \textbf{Debts ($-$)} include money that you \textbf{owe}, a bill you pay or a withdrawal.

\end{itemize}


When we add integers that represent assets and debts, we are finding the \textbf{net balance}.

\[ \textbf{Net Balance} = \textrm{Assets} + \textrm{Debts} \]

\begin{remark}
A statement such as ``a debt of \$50'' is represented mathematically as $-50$. 
The negative sign is omitted in the phrase because the word \emph{debt} already implies a negative value.
\end{remark}

\begin{example}  \textcolor{glaucous}{\textbf{Combining Assets}} 
    
    If a child has \$100 ($+100$) in her piggy bank and receives \$30 ($+30$) for doing her chores, her net 
     balance is \textbf{\$130}.
    \[ 100 + 30 = 130 \]
\end{example}

\begin{example} \textcolor{glaucous}{\textbf{Combining Debts}} 
    
    If you have a credit card debt of \$500 ($-500$) and make another credit card purchase of 
    \$100 ($-100$), your net balance is a \textbf{debt of \$600 ($-600$)}.
    \[ -500 + (-100) = -600 \]
\end{example}

\begin{example} \textcolor{glaucous}{\textbf{Partial Cancellation}}  
    
    If you have \$20 ($+20$) and owe a friend \$25 ($-25$).  Your \$20 isn't enough to 
    cover the entire debt, it only cancels \textbf{part} of it.  After you pay your friend \$20, you will still owe \$5, so
    your net balance is a \textbf{debt of \$5 ($-5$)}.
    \[ 20 + (-25) = -5 \]
\end{example}

\begin{example} \textcolor{glaucous}{\textbf{Full Cancellation}} 
    
    If a taxpayer owes the IRS \$300 ($-300$) and has \$800 ($+800$) in the bank, they can pay the 
    entire debt and keep the change.  Their net balance is \textbf{\$500 ($+500$)}.
    \[ -300 + 800 = 500 \]
\end{example}


\subsection*{Key Observations}

If you take a close look at the examples above you will notice the following relationships between the signs of the numbers that were added and the 
operation we actually performed.


\begin{itemize}
    \item[] \textbf{Same Signs:} When we add two assets or two debts, we \textbf{add} their values.  The sign of the answer is the 
    same as the sign of the numbers.

    \item[] \textbf{Different Signs:} When we add an asset and a debt, all or part of the debt will be removed or cancelled, 
    so even though we are technically \emph{adding} integers, the arithmetic performed is \textbf{subtraction}.  The sign of 
    the answer depends on which is greater: the asset or the debt.
\end{itemize}



\begin{problem}
Your checking account is overdrawn by \$50 and the bank charges you a \$15 overdraft fee.  Write the addition problem to find your 
net balance.  Then calculate your net balance. 


\begin{multipleChoice}
\choice{$50 + (-15)$}
\choice{$-50 + 15$}
\choice{$50+15$}
\choice[correct]{$-50 + (-15)$}
\end{multipleChoice}


\textcolor{orange!75!black}{\textbf{Net Balance}} = $\answer{-65}$

\end{problem}


\begin{problem}
A soccer team spent \$25 on supplies for a car wash.  By the end of the day, they earned \$175. Write the addition problem to find 
the team's net balance.  Then calculate their net balance. 


\begin{multipleChoice}
\choice{$25 + 175$}
\choice[correct]{$-25 + 175$}
\choice{$-25+(-175)$}
\choice{$25 + (-175)$}
\end{multipleChoice}


\textcolor{orange!75!black}{\textbf{Net Balance}} = $\answer{150}$


\end{problem}


\begin{problem}
It's your birthday!  You receive \$300 as a gift, but you owe your roommate \$500 for rent. Write the addition problem to 
find your net balance after you use your gift money toward the rent.  Then calculate your net balance.

\begin{multipleChoice}
\choice{$-300 + (-500)$}
\choice{$300 + 500$}
\choice[correct]{$300+(-500)$}
\choice{$-300 + 500$}
\end{multipleChoice}

\textcolor{orange!75!black}{\textbf{Net Balance}} = $\answer{-200}$

\end{problem}

\begin{problem}
Think in terms of assets and debts to find the following sums. 

~ \\[0.5em]

$ 80 + (-2)  \begin{prompt}=\answer{78}\end{prompt} $

$ -9 + (-1)  \begin{prompt}=\answer{-10}\end{prompt} $

$ -30 + 27 \begin{prompt}=\answer{-3}\end{prompt} $

\end{problem}


\end{document}